

\section{The monochromatic field with a Gaussian turn-on}
We study the motion of a charged particle in a plane wave field with a Gaussian turn-on.\\
We define the envelope as
\begin{align}
  f(\chi)= \begin{cases}e^{-\frac{\chi^{2}}{(2 \pi \sigma)^{2}}}, & \chi<\chi_{0}  \label{I.1.3.1}\\ 1, & \chi \geq \chi_{0}\end{cases}
\end{align}
and the vector potential in the Lorentz-Coulomb gauge as
\begin{align}
  \mathbf{A}(\phi)=A_{0} f(k \phi)\left(\mathbf{e}_{x} \zeta_{x} \cos \left(k \phi+\Phi_{0}\right)+\mathbf{e}_{y} \zeta_{y} \sin \left(k \phi+\Phi_{0}\right)\right) . \label{I.1.3.2}
\end{align}
where $k$ is the wave number $k=\frac{\omega}{c}, \sigma$ gives the width of the Gaussian envelope, $\Phi_{0}$ is the initial phase and $\zeta_{x}, \zeta_{y}$ are the polarization amplitudes, obeying the condition $\zeta_{x}^{2}+\zeta_{y}^{2}=1$.
It is convenient to define the unity vector of the field propagation direction as $\mathbf{n}=(0,0,1)$ and the four-vector
\begin{align}
  n=(1, \mathbf{n}) \label{I.1.3.3}
\end{align}
which in the case of the propagation along the $O z$ axis becomes $n=(1,0,0,1)$. Also the notation for the dimensionless field amplitude is
\begin{align}
  \xi_{0}=\frac{e_{0} A_{0}}{m c} . \label{I.1.3.4}
\end{align}

\section{Initial conditions}
We denote the initial moment in the laboratory frame by $t_{0}$ and the initial position of the particle by $\boldsymbol{\mathcal { R }}_{0}$. We assume that at $t_{0}$ the particle velocity is $\mathbf{v}_{0}$; the corresponding four-velocity will be denoted by $u_{0}$, and the four-momentum by $p_{0}$. The field propagates along the $\mathbf{n}$ direction, chosen parallel to $O z$, and is described by the potential $A(\phi)$, where $\phi=c t-z$. We denote by $A_{0}$ the value of the potential $A$ in the initial position of the particle $p m b \mathcal{R}_{0}$ at the initial moment $t_{0}$ i.e.
\begin{align}
  A_{0} \equiv\left(0, \mathbf{A}_{\mathbf{0}}\right)=A\left(\phi_{0}\right), \quad \phi_{0}=c t_{0}-\mathbf{k}_{0} \cdot \mathbf{n} . \label{II.2.2.1}
\end{align}
For the case of a pulse, the most natural choice for $t_{0}$ is a moment sufficiently far in the past, when the pulse has not reached the origin and the particle is free (i.e. $A_{0}=0$ ). The general equations of motion will be written for arbitrary values of $t_{0}$ and $\mathbf{v}_{0}$,particular results will be presented for two cases of interest:
\begin{enumerate}
  \item Gaussian pulse; the particle is at rest in $\boldsymbol{\mathcal { R }}_{0}$ at a moment $t_{0}$ very far in the past, when the pulse has not yet reached the point $\boldsymbol{k}_{0}$.
  \item monochromatic field; at the moment $t_{0}$ the particle is located in $\mathbf{k}_{0}$, with an initial velocity $\mathbf{v}_{0}$, the value of the vector potential in the initial position of the particle being $A_{0}=A\left(c t_{0}-\mathbf{k}_{0} \cdot \mathbf{n}\right)$.
  \item We will combine the two cases into a general one, corresponding to a field with a gaussian turn-on, followed by a monochromatic oscillation of infinite length. The initial conditions are the same as in the previous case (1).
\end{enumerate}

\section[The trajectory for arbitrary initial conditions]{The trajectory for arbitrary initial conditions}
In the relativistic case the equations of motion are
\begin{align}
  \frac{d p^{\mu}}{d \tau}=e_{0} F^{\mu v} u_{v}, \quad u=\frac{d x}{d \tau}, \quad p=m_{0} u, \label{II.2.3.1}
\end{align}
where $m_{0}, e_{0}$ are the rest mass and charge of the particle, $\tau$ is the proper time and $\hat{\mathbf{F}}$ is the electromagnetic four-tensor of the field intensities:
\begin{align}
  F^{\mu v}=\partial^{\mu} A^{v}-\partial^{v} A^{\mu} \equiv \left\lvert\, \begin{array}{cccc}
                                                                             0                             & \frac{d A^{1}(\phi)}{d \phi} & \frac{d A^{2}(\phi)}{d \phi} & 0                             \\
                                                                             -\frac{d A^{1}(\phi)}{d \phi} & 0                            & 0                            & -\frac{d A^{1}(\phi)}{d \phi} \\
                                                                             -\frac{d A^{2}(\phi)}{d \phi} & 0                            & 0                            & -\frac{d A^{2}(\phi)}{d \phi} \\
                                                                             0                             & \frac{d A^{1}(\phi)}{d \phi} & \frac{d A^{2}(\phi)}{d \phi} & 0
                                                                           \end{array}\right\rvert . \label{II.2.3.2}
\end{align}
With the above definitions the equations of motion become
\begin{align}
   & \frac{d p^{0}}{d \tau}=\frac{d p^{3}}{d \tau}=-e_{0} \frac{d \mathbf{A}(\phi)}{d \phi} \cdot \mathbf{u}_{\perp},  \label{II.2.3.3} \\
   & \frac{d \mathbf{p}_{\perp}}{d \tau}=-e_{0} \frac{d \mathbf{A}(\phi)}{d \phi}\left(u^{0}-u^{3}\right) . \label{II.2.3.4}            \\
\end{align}
From the first of the above formulae one obtains
\begin{align}
  \frac{d\left(p^{0}-p^{3}\right)}{d \tau}=0 \label{II.2.3.5}
\end{align}
i.e. $p^{0}-p^{3}=m_{0}\left(u^{0}-u^{3}\right)$ is a constant of motion; taking into account the initial conditions we have $p^{0}-p^{3}=m_{0}\left(u_{0}^{0}-u_{0}^{3}\right)$. Further, noticing that
\begin{align}
  u^{0}-u^{3}=\text { const }=\frac{d}{d \tau}\left(x^{0}-x^{3}\right)=\frac{d \phi}{d \tau} \label{II.2.3.6}
\end{align}
we have
\begin{align}
  \frac{d \phi}{d \tau}=\left(u_{0}^{0}-u_{0}^{3}\right)=\frac{n \cdot p_{0}}{m_{0}} \quad \longrightarrow \quad \phi(\tau)=\frac{n \cdot p_{0}}{m_{0}}\left(\tau-\tau_{0}\right)+\phi_{0} . \label{II.2.3.7}
\end{align}
Eq.(\ref{II.2.3.4}) becomes
\begin{align}
  \frac{d \mathbf{p}_{\perp}}{d \tau}=-e_{0} \frac{d \mathbf{A}}{d \phi} \frac{d \phi}{d \tau} . \label{II.2.3.8}
\end{align}
It is again convenient to look for the solutions of the above equations as functions of $\phi$. Using the initial conditions one obtains
\begin{align}
  \mathbf{p}_{\perp}(\phi)=\left(-e_{0} \mathbf{A}(\phi)+e_{0} \mathbf{A}_{0}\right)+\mathbf{p}_{0 \perp} \label{II.2.3.9}
\end{align}
and, then, from Eqs. (\ref{II.2.3.3}) and (\ref{II.2.3.9})
\begin{align}
   & \frac{d p^{3}(\phi)}{d \phi}=\frac{d p^{3}}{d \tau} \frac{d \tau}{d \phi}=\frac{d p^{3}}{d \tau} \frac{1}{n \cdot u_{0}}=-\frac{e_{0}}{n \cdot u_{0}} \frac{d \mathbf{A}(\phi)}{d \phi} \cdot \mathbf{u}_{\perp}=-\frac{e_{0}}{m_{0} n \cdot u_{0}} \frac{d \mathbf{A}(\phi)}{d \phi}\left(-e_{0} \mathbf{A}(\phi)+e_{0} \mathbf{A}_{0}+\mathbf{p}_{0 \perp}\right)  \label{II.2.3.10} \\
   & \frac{d p^{3}(\phi)}{d \phi}=\frac{1}{2 n \cdot p_{0}} \frac{d}{d \phi}\left(e_{0} \mathbf{A}(\phi)-e_{0} \mathbf{A}_{0}-\mathbf{p}_{0 \perp}\right)^{2}=\frac{1}{2 n \cdot p_{0}} \frac{d}{d \phi}\left(e_{0}^{2} \boldsymbol{\mathcal { A }}(\phi)\right)^{2} \label{II.2.3.11}                                                                                                      \\
\end{align}
with
\begin{align}
  e_{0} \mathcal{A}(\phi)=\left(e_{0} \mathbf{A}(\phi)-e_{0} \mathbf{A}_{0}+\mathbf{p}_{0 \perp}\right) . \label{II.2.3.12}
\end{align}
from which we get the solution
\begin{align}
  p^{3}(\phi)=\frac{e_{0}^{2} \mathcal{A}^{2}(\phi)}{2 n \cdot p_{0}}-\frac{\mathbf{p}_{\perp 0}^{2}}{2 n \cdot p_{0}}+p_{0}^{3}, \label{II.2.3.13}
\end{align}
An equivalent notation for the above formula is
\begin{align}
  e_{0} \mathcal{A}\left(\phi_{0}\right)=e_{0} \mathcal{A}_{0}=\mathbf{p}_{0 \perp}, \quad p^{3}(\phi)=\frac{e_{0}^{2} \mathcal{A}^{2}(\phi)}{2 n \cdot p_{0}}-\frac{e_{0}^{2} \mathcal{A}_{0}^{2}}{2 n \cdot p_{0}}+p_{0}^{3} . \label{II.2.3.14}
\end{align}
Finally, the solution for the four-momentum of the particle is
\begin{align}
   & \mathbf{p}_{\perp}(\phi)=-e_{0} \boldsymbol{\mathcal { A }}(\phi)  \label{II.2.3.15}                                                                    \\
   & p^{3}(\phi)=\frac{e_{0}^{2} \mathcal{A}^{2}(\phi)}{2 n \cdot p_{0}}-\frac{e_{0}^{2} \mathcal{A}_{0}^{2}}{2 n \cdot p_{0}}+p_{0}^{3},  \label{II.2.3.16} \\
   & p^{0}(\phi)=p^{3}(\phi)+n \cdot p_{0} . \label{II.2.3.17}                                                                                               \\
\end{align}
The differential equations in $\phi$ obeyed by the coordinates are
\begin{align}
  \frac{d \mathbf{r}}{d \phi}=\frac{d \mathbf{r}}{d \tau} \frac{d \tau}{d \phi}=\frac{1}{\left(u_{0}^{0}-u_{0}^{3}\right)} \frac{d \mathbf{r}}{d \tau}=\frac{\mathbf{p}(\phi)}{n \cdot p_{0}} . \label{II.2.3.18}
\end{align}
with the solution
\begin{align}
   & \mathbf{r}_{\perp}=\boldsymbol{R}_{0 \perp}-\frac{1}{n \cdot p_{0}} \int_{\phi_{0}}^{c t-z} e_{0} \boldsymbol{\mathcal { A }}(\chi) d \chi  \label{II.2.3.19}                                                                                            \\
   & z=\mathcal{R}_{0}^{3}+\int_{\phi_{0}}^{c t-z} d \chi\left[\frac{e_{0}^{2} \boldsymbol{\mathcal { A }}^{2}(\chi)-e_{0}^{2} \boldsymbol{\mathcal { A }}_{0}^{2}}{2\left(n \cdot p_{0}\right)^{2}}+\frac{p_{0}^{3}}{n \cdot p_{0}}\right] \label{II.2.3.20} \\
\end{align}


\subsection{The case of a particle initially at rest}
The motion along the turn-on of the pulse\\
We assume that the initial positions of the particles are all in the plane $O x y$, i.e. $\mathbf{R}_{0} \cdot \mathbf{n}=0$. We chose $t_{0} \ll-\sigma$, so that the initial value of the vector potential is $A_{0}=(0,0,0,0)$; the initial value of $\tau$ (the proper time) will be conventionally taken as $\tau_{0}=t_{0}$. Since the particle is initially at rest, the initial four-velocity is $u_{0}=(c, 0,0,0)$ and the initial four-momentum is $p_{0}=(m c, 0,0,0)$; using (\ref{II.2.3.6}) we can write
\begin{align}
  \frac{d \phi}{d \tau}=c . \label{II.2.4.1}
\end{align}
i.e. we can express $\phi$ as a function of $\tau$ : $\phi(\tau)=c\left(\tau-\tau_{0}\right)+\phi_{0}$, with $\phi_{0}=c t_{0}=c \tau_{0}$, and finally
\begin{align}
  \phi(\tau)=c \tau \label{II.2.4.2}
\end{align}
with these we can calculate explicitly the 4 -momentum $p$ of the particle as a function of $\tau$ :
\begin{align}
  p^{3}(\tau)=\frac{e^{2} \mathbf{A}^{2}(c \tau)}{2 m c}, \quad p_{\perp}(\tau)=-e \mathbf{A}(c \tau) . \label{II.2.4.3}
\end{align}
The trajectory of the particle is given by
\begin{align}
  \mathbf{r}_{\perp}(\tau)=\mathbf{k}_{0 \perp}-\frac{e}{m c} \int_{-\infty}^{c \tau} \mathbf{A}(\chi) d \chi, \quad z(\tau)=\mathcal{K}_{03}+\frac{e^{2}}{2(m c)^{2}} \int_{-\infty}^{c \tau} A^{2}(\chi) d \chi \label{II.2.4.4}
\end{align}
The motion in the monochromatic part of the field\\
Now, we can study the motion in the monochromatic part of the field, i.e. for $\phi \in(0,+\infty)$, using as initial conditions $\mathbf{r}(0)$, $\mathbf{P}(0)$,calculated at the end of the turn-on of the pulse. The field on the turn=on is
\begin{align}
  \mathbf{A}(\phi)=A_{0} e^{-\frac{\phi^{2}}{(2 \pi \sigma)^{2}}}\left(\mathbf{e}_{x} \zeta_{x} \cos \left(k \phi+\Phi_{0}\right)+\mathbf{e}_{y} \zeta_{y} \sin \left(k \phi+\Phi_{0}\right)\right) . \label{II.2.4.5}
\end{align}
and we have the initial conditions on the monochromatic part of the field
\begin{align}
   & \phi(0)=0  \label{II.2.4.6}                                                                                                                                                                                                                                                                                                                                                                      \\
   & \mathbf{p}_{\perp}(0)=-e_{0} \mathbf{A}(0)=-e_{0} A_{0}\left(\mathbf{e}_{x} \zeta_{x} \cos \left(\Phi_{0}\right)+\mathbf{e}_{y} \zeta_{y} \sin \left(\Phi_{0}\right)\right), \quad p^{3}(0)=\frac{e_{0}^{2} \mathbf{A}(0)^{2}}{2 m c}=\frac{e_{0}^{2} A_{0}^{2}\left(\zeta_{x}^{2} \cos ^{2}\left(\Phi_{0}\right)+\zeta_{y}^{2} \sin ^{2}\left(\Phi_{0}\right)\right)}{2 m c} . \label{II.2.4.7} \\
\end{align}
For the initial position one can prove that is $\sigma$ is large enough, at $t=0$ the transverse motion is negligible, i.e.
\begin{align}
  \mathbf{r}_{\perp}(0)=\mathbf{R}_{0 \perp}-\frac{e}{m c} \int_{-\infty}^{0} \mathbf{A}(\chi) d \chi \approx \mathbf{k}_{0 \perp} \label{II.2.4.8}
\end{align}
Note that $z(0)$ is not zero, only $\phi(0)$ is zero, so the initial position of the particle is not in the plane $O x y$ but in a plane parallel to it, at a distance $z(0)$ from it. We have
\begin{align}
  z(0)=\mathcal{K}_{0 z}+\frac{e^{2}}{2(m c)^{2}} \int_{-\infty}^{0} A^{2}(\chi) d \chi=\mathcal{K}_{0 z}+\frac{e^{2} A_{0}^{2}}{2(m c)^{2}} \int_{-\infty}^{0} e^{-\frac{\chi^{2}}{2(\pi \sigma)^{2}}}\left(\zeta_{x}^{2} \cos ^{2}\left(k \chi+\Phi_{0}\right)+\zeta_{y}^{2} \sin ^{2}\left(k \chi+\Phi_{0}\right)\right) d \chi \label{II.2.4.9}
\end{align}
Again, in the limit of large $\sigma$ we have
\begin{align}
  z(0) \approx \mathcal{K}_{0 z}+\frac{e^{2} A_{0}^{2}}{2(m c)^{2}} \frac{\pi^{3 / 2} \sigma}{2 \sqrt{2} k}=\mathcal{R}_{0 z}+\xi_{0}^{2} \frac{\pi^{3 / 2} \sigma}{4 \sqrt{2} k} . \label{II.2.4.10}
\end{align}
Next, we apply the general formulae for the trajectory in the monochromatic part of the field, using the above initial conditions.
\begin{align}
  p_{0} \rightarrow p(0), \quad A_{0} \rightarrow A(0), \quad \mathbf{r}(0) \rightarrow \mathbf{r}(0), \quad \phi_{0} \rightarrow \phi(0)=0 \label{II.2.4.11}
\end{align}
We calculate the component 0 of the momentum as
\begin{align}
  p^{0}(0)=\sqrt{\left(m_{0} c\right)^{2}+\mathbf{p}^{2}}=\sqrt{\left(m_{0} c\right)^{2}+\mathbf{p}_{\perp}^{2}+\left(p^{3}\right)^{2}}=\sqrt{\left(m_{0} c\right)^{2}+e_{0}^{2} \mathbf{A}^{2}(0)+\left(\frac{e_{0}^{2} \mathbf{A}^{2}(0)}{2 m c}\right)^{2}}=m_{0} c+\frac{e_{0}^{2} \mathbf{A}^{2}(0)}{2 m c} \label{II.2.4.12}
\end{align}
The previous result is in agreement with the general conservation law for $n \cdot p_{0}$, i.e.
\begin{align}
  n \cdot p_{0}=p^{0}(0)-p^{3}(0)=m_{0} c . \label{II.2.4.13}
\end{align}
with the new initial conditions we have the momentum during the flat part of the pulse
\begin{align}
   & \mathbf{p}_{\perp}(\phi)=-e_{0} \mathcal{A}(\phi)=-e_{0} \mathbf{A}(\phi)  \label{II.2.4.14}                                                                                                                                                                                                                                                                                                                                       \\
   & p^{3}(\phi)=\frac{e_{0}^{2} \mathcal{A}^{2}(\phi)}{2 n \cdot p_{0}}-\frac{e_{0}^{2} \mathcal{A}_{0}^{2}}{2 n \cdot p_{0}}+p_{0}^{3}=\frac{e_{0}^{2} \mathbf{A}^{2}(\phi)}{2 m c}-\frac{\mathbf{p}_{\perp}^{2}(0)}{2 m c}+p^{3}(0)=\frac{e_{0}^{2} \mathbf{A}^{2}(\phi)}{2 m c}-\frac{e_{0}^{2} \mathbf{A}^{2}(0)}{2 m c}+\frac{e_{0}^{2} \mathbf{A}^{2}(0)}{2 m c}=\frac{e_{0}^{2} \mathbf{A}^{2}(\phi)}{2 m c}  \label{II.2.4.15} \\
   & p^{0}(\phi)=m c+p^{3}(\phi) \label{II.2.4.16}                                                                                                                                                                                                                                                                                                                                                                                      \\
\end{align}
which now can be calculated explicitly
\begin{align}
   & \mathbf{p}_{\perp}(\phi)=-e_{0} A_{0}\left(\mathbf{e}_{x} \zeta_{x} \cos \left(k \phi+\Phi_{0}\right)+\mathbf{e}_{y} \zeta_{y} \sin \left(k \phi+\Phi_{0}\right)\right),  \label{II.2.4.17} \\
   & p^{3}(\phi)=\frac{e_{0}^{2} A_{0}^{2}}{2 m c}\left(\zeta_{x}^{2} \cos ^{2}\left(k \phi+\Phi_{0}\right)+\zeta_{y}^{2} \sin ^{2}\left(k \phi+\Phi_{0}\right)\right) . \label{II.2.4.18}       \\
\end{align}
or, as a function of tau, with $\phi=c \tau$,
\begin{align}
   & \mathbf{p}_{\perp}(\tau)=-e_{0} A_{0}\left(\mathbf{e}_{x} \zeta_{x} \cos \left(k c \tau+\Phi_{0}\right)+\mathbf{e}_{y} \zeta_{y} \sin \left(k c \tau+\Phi_{0}\right)\right),  \label{II.2.4.19} \\
   & p^{3}(\tau)=\frac{e_{0}^{2} A_{0}^{2}}{2 m c}\left(\zeta_{x}^{2} \cos ^{2}\left(k c \tau+\Phi_{0}\right)+\zeta_{y}^{2} \sin ^{2}\left(k c \tau+\Phi_{0}\right)\right) . \label{II.2.4.20}       \\
\end{align}
Now we can write the trajectory explicitly as a function of $\tau$ :
\begin{align}
  \mathbf{r}_{\perp}(\tau) & =\mathbf{R}_{0 \perp}+\int_{0}^{\tau} d \tau^{\prime} \frac{\mathbf{p}_{\perp}\left(\tau^{\prime}\right)}{m}=\mathbf{R}_{0 \perp}-\frac{e A_{0}}{m} \int_{0}^{\tau} d \tau^{\prime}\left(\mathbf{e}_{x} \zeta_{x} \cos \left(k c \tau^{\prime}+\Phi_{0}\right)+\mathbf{e}_{y} \zeta_{y} \sin \left(k c \tau^{\prime}+\Phi_{0}\right)\right)  \label{II.2.4.21} \\
                           & =\mathbf{R}_{0 \perp}-\frac{e A_{0}}{m}\left[\mathbf{e}_{x} \zeta_{x} \frac{\sin \left(k c \tau+\Phi_{0}\right)-\sin \left(\Phi_{0}\right)}{k c}-\mathbf{e}_{y} \zeta_{y} \frac{\cos \left(k c \tau+\Phi_{0}\right)-\cos \left(\Phi_{0}\right)}{k c}\right]                                                                                                    \\
                           & =\mathbf{R}_{0 \perp}-\frac{\xi_{0}}{k}\left[\mathbf{e}_{x} \zeta_{x}\left(\sin \left(k c \tau+\Phi_{0}\right)-\sin \left(\Phi_{0}\right)\right)-\mathbf{e}_{y} \zeta_{y}\left(\cos \left(k c \tau+\Phi_{0}\right)-\cos \left(\Phi_{0}\right)\right)\right]                                                                                                    \\
                           & =\mathbf{R}_{0 \perp}+\frac{\xi_{0}}{k}\left[\mathbf{e}_{x} \zeta_{x} \sin \left(\Phi_{0}\right)-\mathbf{e}_{y} \zeta_{y} \cos \left(\Phi_{0}\right)\right]-\frac{\xi_{0}}{k}\left[\mathbf{e}_{x} \zeta_{x} \sin \left(k c \tau+\Phi_{0}\right)-\mathbf{e}_{y} \zeta_{y} \cos \left(k c \tau+\Phi_{0}\right)\right]  \label{II.2.4.22}                         \\
                           & =\mathbf{r}_{i \perp}+\delta \mathbf{r}_{\perp}(t) \label{II.2.4.23}                                                                                                                                                                                                                                                                                           \\
\end{align}
\begin{align}
  z(\tau) & =z(0)+\int_{0}^{\tau} d \tau^{\prime} \frac{p^{3}\left(\tau^{\prime}\right)}{m}=z(0)+\frac{e^{2} A_{0}^{2}}{2 m^{2} c} \int_{0}^{\tau} d \tau^{\prime}\left(\zeta_{x}^{2} \cos ^{2}\left(k c \tau^{\prime}+\Phi_{0}\right)+\zeta_{y}^{2} \sin ^{2}\left(k c \tau^{\prime}+\Phi_{0}\right)\right) \\
          & =\mathcal{R}_{0 z}+\xi_{0}^{2} \frac{\pi^{3 / 2} \sigma}{4 \sqrt{2} k}+\frac{\xi_{0}^{2}}{4}\left[c \tau+\frac{\sin \left(2 k c \tau+2 \Phi_{0}\right)-\sin \left(2 \Phi_{0}\right)}{2 k}\right]                                                                                                 \\
          & =\mathfrak{R}_{0 z}+\frac{\xi_{0}^{2}}{8 k}\left[\sqrt{2} \pi^{3 / 2} \sigma-\sin \left(2 \Phi_{0}\right)\right]+\frac{\xi_{0}^{2}}{4} c \tau+\frac{\xi_{0}^{2}}{8 k} \sin \left(2 k c \tau+2 \Phi_{0}\right)                                                                                    \\
          & =z_{i}+\delta z(t) \label{II.2.4.24}
\end{align}
The trajectory as a function of $\phi$ is obtained by the substitution $\phi=c \tau$ :
\begin{align}
   & \mathbf{r}_{\perp}(\phi)=\mathbf{R}_{0 \perp}+\frac{\xi_{0}}{k}\left[\mathbf{e}_{x} \zeta_{x} \sin \left(\Phi_{0}\right)-\mathbf{e}_{y} \zeta_{y} \cos \left(\Phi_{0}\right)\right]-\frac{\xi_{0}}{k}\left[\mathbf{e}_{x} \zeta_{x} \sin \left(k \phi+\Phi_{0}\right)-\mathbf{e}_{y} \zeta_{y} \cos \left(k \phi+\Phi_{0}\right)\right],  \label{II.2.4.25} \\
   & z(\phi)=\mathcal{R}_{0 z}+\frac{\xi_{0}^{2}}{8 k}\left[\sqrt{2} \pi^{3 / 2} \sigma-\sin \left(2 \Phi_{0}\right)\right]+\frac{\xi_{0}^{2}}{4} \phi+\frac{\xi_{0}^{2}}{8 k} \sin \left(2 k \phi+2 \Phi_{0}\right) . \label{II.2.4.26}                                                                                                                         \\
\end{align}


\subsection{The case of a particle with an initial momentum along $O z$}
The motion along the turn-on of the pulse\\
We consider the same field as in the previous section, but now we assume that the particle has an initial momentum $p_{0} \equiv\left(p_{0}^{0}, 0,0, p_{0}^{3}\right)$ at a moment $t_{0}$ sufficiently far in the past, when the pulse has not yet reached the point $\boldsymbol{K}_{0}$. The initial value of the proper time is $\tau_{i}=-t_{0}$ and the initial conditions are
\begin{align}
  \mathbf{r}_{0}=\mathbf{k}_{0} \quad \phi\left(\tau_{i}\right)=c t_{0}-\mathcal{R}_{0 z}=c t_{0}, \quad \mathbf{p}_{\perp}(0)=\mathbf{0}, \quad p^{3}(0)=p_{0}^{3}, \quad p^{0}(0)=\sqrt{\left(m_{0} c\right)^{2}+\left(p_{0}^{3}\right)^{2}} . \label{II.2.5.1}
\end{align}
Now, the dependence of $\phi$ on $\tau$ is
\begin{align}
  \frac{d \phi}{d \tau}=\frac{n \cdot p_{0}}{m_{0}} \quad \longrightarrow \quad \phi(\tau)=\frac{n \cdot p_{0}}{m_{0}}\left(\tau-\tau_{i}\right)+\phi\left(\tau_{i}\right)=\frac{n \cdot p_{0}}{m_{0}}\left(\tau-t_{0}\right)+c t_{0} . \label{II.2.5.2}
\end{align}
The moment  $\tau$ when the turn-on of the pulse is completed, i.e. when $\phi(\tau)=0$, is
\begin{align}
  \tau_{f}=t_{0}-\frac{m_{0}}{n \cdot p_{0}} c t_{0}=t_{0}\left(1-\frac{m_{0} c}{n \cdot p_{0}}\right) . \label{II.2.5.3}
\end{align}
The modified vector potential is
\begin{align}
  e_{0} \mathcal{A}(\phi)=e_{0} \mathbf{A}(\phi)-e_{0} \mathbf{A}_{0}+\mathbf{p}_{0 \perp}=e_{0} \mathbf{A}(\phi) \label{II.2.5.4}
\end{align}
The four-momentum of the particle during the turn-on of the pulse is given by
\begin{align}
  \mathbf{p}_{\perp}(\phi)=-e_{0} \mathbf{A}(\phi), \quad p^{3}(\phi)=\frac{e_{0}^{2} \mathbf{A}^{2}(\phi)}{2 n \cdot p_{0}}+p_{0}^{3}, \quad p^{0}(\phi)=p^{3}(\phi)+n \cdot p_{0}=\frac{e_{0}^{2} \mathbf{A}^{2}(\phi)}{2 n \cdot p_{0}}+p_{0}^{0} \label{II.2.5.5}
\end{align}

The momentum of the particle at $\phi=0$ is
\begin{align}
  {\bf p}_\perp(0) = -e_0{\bf A}(0), \quad p^{3}(0)=\frac{e_{0}^{2} \mathbf{A}^{2}(0)}{2 n \cdot p_{0}}+p_{0}^{3}, \quad p^{0}(0)=p^{3}(0)+n \cdot p_{0}=\frac{e_{0}^{2} \mathbf{A}^{2}(0)}{2 n \cdot p_{0}}+p_{0}^{0}
\end{align}
This value will serve as initial momentum for the remaining of the trajectory, denoted by $\widetilde p_0$
\begin{align}
  \widetilde p_0=(p_0^0 + \frac{e_{0}^{2} \mathbf{A}^{2}(0)}{2 n \cdot p_{0}},-e_0A_x(0), -e_0A_y(0), p_0^3 + \frac{e_{0}^{2} \mathbf{A}^{2}(0)}{2 n \cdot p_{0}})
\end{align}
For the vector potential (\ref{I.1.3.2}) we have
\begin{align}
  {\bf A}(0) = A_0({\bf e}_x \zeta_x\cos(\Phi_0)+{\bf e}_y \zeta_y\sin(\Phi_0))
\end{align}
and
\begin{align}
  \widetilde p_0=(p_0^0 + \frac{e_{0}^{2} \mathbf{A}^{2}(0)}{2 n \cdot p_{0}},-e_0A_x(0), -e_0A_y(0), p_0^3 + \frac{e_{0}^{2} \mathbf{A}^{2}(0)}{2 n \cdot p_{0}})
\end{align}

The position of the particle at $\phi=0$ will be, using (\ref{II.2.3.19}, \ref{II.2.3.20})
\begin{align}
   & \mathbf{r}_{\perp}(0)=\boldsymbol{R}_{0 \perp}-\frac{1}{n \cdot p_{0}} \int_{-\infty}^{0} e_{0} \boldsymbol{\mathcal { A }}(\chi) d \chi  \sim \boldsymbol{R}_{0 \perp}                                                                \\
   & z=\mathcal{R}_{0}^{3}+\int_{\phi_{0}}^{c t-z} d \chi\left[\frac{e_{0}^{2} \boldsymbol{\mathcal { A }}^{2}(\chi)-e_{0}^{2} \boldsymbol{\mathcal { A }}_{0}^{2}}{2\left(n \cdot p_{0}\right)^{2}}+\frac{p_{0}^{3}}{n \cdot p_{0}}\right] \\
\end{align}
